%% kapa: exam=1 round=36 year=2025 period=1 form=A subject=civil_law subject_ko=민법 seq=33 number=33
%% source: Q-Net 감정평가사 기출문제 (cst00309) artlSeq=5231525 file=raw/1008/5231525/2025년 제36회 감정평가사 1차시험 1교시 문제지 원본.pdf
%% conversion: kordoc 4.24.51 -> tools/gen1.py (auto); review=proofread-sonnet; flags=none
\begin{question}{33}
지역권에 관한 설명으로 옳지 않은 것은? (다툼이 있으면 판례에 따름)
\begin{choices}
  \item 요역지소유권과 지역권이 함께 이전되지 않도록 하는 약정은 유효하며 이를 등기할 수 있다.
  \item 승역지소유자가 지역권자의 지역권 행사를 위하여 공작물 수선의무를 부담하기로 한 경우, 승역지소유자의 특별승계인도 그 의무를 부담한다.
  \item 동일한 승역지 위에 수개의 용수지역권이 설정될 수 있다.
  \item 전망지역권은 소멸시효에 걸리지 않는다.
  \item 통행지역권을 시효로 취득한 자는 특별한 사정이 없는 한 도로 설치 및 사용에 의하여 승역지소유자가 입은 손해를 보상해야 한다.
\end{choices}
\end{question}
\answer{4}
